\chapter{Local Stores and the Transfer Unit}\label{ch:memory}

\section{The Local Stores}

\subsection{Data Local Store}
The D-LS is a byte-addressed memory private to its core, \texttt{DLS\_SIZE} bytes long
(Section~\ref{sec:lssize}; 64~KiB by default). Loads and stores address it with \reg{rs1} plus a
sign-extended 16-bit offset. An access of size $n$ (1, 2, 4 or 8 bytes) must be naturally aligned and
lie entirely inside $[0, \texttt{DLS\_SIZE})$; otherwise it traps with cause \cause{lsfault} and
\csr{tval} set to the address. Addresses are full 64-bit values: one at or past the size is out of
range whatever its low bits, so no address wraps around into a smaller store. There is no address
translation: D-LS addresses are physical, and address 0 is a valid location.

Because a 16-bit signed offset from \reg{x0} reaches addresses 0--32767, data placed in the lowest
32~KiB of the D-LS (all of a D-LS of 32~KiB or less) can be addressed absolutely, as in
\insn{ld t1, counter(zero)}, without a base register. Trap handlers use this to save registers before
they have any free register.

\subsection{Instruction Local Store}
The I-LS is a memory of \texttt{ILS\_SIZE} bytes (64~KiB by default) from which, and only from which,
instructions are fetched. Before each instruction the core checks that \reg{pc} is a multiple of 4 and
less than \texttt{ILS\_SIZE}; if not, the fetch traps with cause \cause{ifault} and \csr{tval} =
\reg{pc}. No load or store can read or write the I-LS. It is filled at reset from the boot image and at
run time only by \insn{tgetx}.

\subsection{The Local Stores' sizes}\label{sec:lssize}
\texttt{DLS\_SIZE} and \texttt{ILS\_SIZE} are properties of the platform: each a power of two from
4~KiB to 64~KiB, chosen independently, the same on every core of a machine, and fixed while it runs.
The platform table's header gives both in KiB (Section~\ref{sec:stpd}). The smallest, 4~KiB, holds the
boot ROM that a core-control reset reloads into each store (Section~\ref{sec:corectl}). Every check that
names a Local Store uses its size: loads, stores, \insn{ltag} and \insn{linst} (the set's entry) and
the D-LS range of \insn{tget} and \insn{tput} against \texttt{DLS\_SIZE}; the fetch, the I-LS range of
\insn{tgetx} and the I-LS address of \insn{xmap} against \texttt{ILS\_SIZE}. A smaller store adds no
fault: an address the store does not have traps exactly as an address past 64~KiB does.

Software learns the sizes from the platform table; system software reads it at boot. Every core also
starts with the D-LS size in \reg{sp} (Section~\ref{sec:reset}), so that a bare-metal program, on any
core and without a device capability, has its stack at the top of its D-LS without knowing the number.
A program image whose Local Store sections do not fit the machine's sizes cannot be loaded; the
simulators refuse it.

\section{Far spaces}

Everything beyond the Local Stores is organised into \term{spaces}, each a separate address range:

\begin{center}\small
\begin{tabular}{@{}c l l l@{}}
\toprule
\textbf{Space} & \textbf{Name} & \textbf{Granule} & \textbf{Notes}\\\midrule
0 & none & --- & the space of the null capability\\
2 & T2 cluster store & 8 bytes & memory\\
3 & T3 system memory & 8 bytes & memory\\
4 & T4 device space & 1 byte & ports, registers and read-only windows (Chapter~\ref{ch:platform})\\
5 & core control & 8 bytes & one doubleword per core; \insn{tput} only (Section~\ref{sec:corectl})\\
6 & endpoints & --- & names of nodes for \insn{csend}/\insn{crecv}; no transfers (Section~\ref{sec:endpoints})\\
\bottomrule
\end{tabular}
\end{center}

A far location is never named by a bare address. Every far access gives a capability register and an
offset; the access goes to \texttt{base + offset} in the capability's space and must lie inside the
capability's segment. The base of every capability is a multiple of 8 in every space
(Section~\ref{sec:capalign}), so the granule of a space constrains only the offsets and lengths of
accesses.

\section{Transfers}\label{sec:transfers}

A \term{transfer} copies a block of bytes between a Local Store and a far space:

\begin{center}\small
\begin{tabular}{@{}l l l l@{}}
\toprule
\textbf{Instruction} & \textbf{Direction} & \textbf{Permission} & \textbf{Local Store}\\\midrule
\insref{tget} \texttt{tag, rl, cs, ro, rn} & far $\to$ local & \perm{R} & D-LS\\
\insref{tput} \texttt{tag, rl, cs, ro, rn} & local $\to$ far & \perm{W} & D-LS\\
\insref{tgetx} \texttt{tag, rl, cs, ro, rn} & far $\to$ local & \perm{X} & I-LS (memory spaces only)\\
\bottomrule
\end{tabular}
\end{center}

\reg{rl} holds the Local Store address, \reg{ro} the offset in the segment and \reg{rn} the length in
bytes. A transfer is checked in full when it issues; if any check fails, the instruction traps
precisely, before any byte has moved. The checks, in order, are: the capability is live
(\cause{capinv}); it has the required permission and names a space the transfer may use: never an
endpoint, \insn{tgetx} never device or core-control space, and core-control space only for \insn{tput}
(\cause{perm}); outside device space, \reg{rl}, \reg{ro} and \reg{rn} are multiples of 8
(\cause{xalign}); the far range lies inside the segment (\cause{bounds}); and the Local Store range lies
inside the Local Store (\cause{lsfault}). For every transfer fault \csr{tval} is the far offset
\reg{ro}, except \cause{lsfault}, for which it is the Local Store address \reg{rl}.

A \insn{tget} or \insn{tput} that has passed its checks is then ordered after the core's pending
device reads: if its D-LS range overlaps the destination of a device read that is still in flight, it
waits at issue until that read has completed (Section~\ref{sec:devxfer}). Last, it waits for its tag.

Once it has issued, a memory transfer always completes; there are no asynchronous transfer errors. A
device read completes when its device has the data, which for an event port can be arbitrarily late.

\section{As-if-synchronous semantics}\label{sec:asif}

\begin{rule*}{The as-if-synchronous rule}
Every transfer takes effect at its issue point in program order: the program behaves exactly as if
the whole block had been copied, instantaneously, when the transfer instruction executed.
\end{rule*}

An implementation is free to perform the copy later, in the background, provided no instruction can
tell the difference. It achieves that by stalling every access that could observe an incomplete
transfer (Table~\ref{tab:hazards}), and by performing the far-side effects of one core's transfers in
issue order.

\begin{table}[htbp]
\centering\small
\caption{Local Store hazards and the stalls that hide them}\label{tab:hazards}
\begin{tabularx}{\linewidth}{@{}l X X@{}}
\toprule
\textbf{In-flight transfer} & \textbf{Its Local Store range is} & \textbf{Accesses that stall while it is in flight}\\\midrule
\insn{tget} & being written by the TU & loads, stores, \insn{ltag}, \insn{linst} that overlap it\\
\insn{tput} & being read by the TU & stores and \insn{linst} that overlap it\\
\insn{tgetx} & being written by the TU (I-LS) & instruction fetch from an overlapping address\\
pending device read & to be written when the device completes the read & loads and stores that overlap
  it; the issue of a \insn{tget} or \insn{tput} whose D-LS range overlaps it\\
\bottomrule
\end{tabularx}
\end{table}

Two consequences matter in practice:
\begin{itemize}
\item A program may load from a buffer immediately after issuing the \insn{tget} that fills it, and may
  refill a buffer immediately after issuing the \insn{tput} that empties it. No \insn{twait} is ever
  needed for correctness.
\item A program may jump to code immediately after issuing the \insn{tgetx} that loads it. There is no
  instruction-synchronisation barrier.
\item A program may start a device read and use its destination at once: a load from it, or a
  \insn{tput} that sends it on, waits for the device. The wait can be long (the next key, the next
  deadline), which is exactly how a program sleeps until an event.
\end{itemize}

Overlapping a transfer with computation (double buffering, prefetching) is therefore purely a
performance choice. The \texttt{stream\_sum} example (Section~\ref{sec:ex-stream}) computes the same
sum with and without double buffering.

\subsection{Waiting explicitly}
\insn{twait} waits for a set of tags, and the pseudo-instruction \insn{tfence} waits for all of them.
\insn{tpoll} returns the busy subset of a set of tags without waiting, and \insn{tany} waits until any
of a set of tags is idle or a gift arrives. These matter only for timing: for example, before reading
\csr{cycle} to time a loop, to overlap work with a transfer whose progress a program wants to observe,
or to wait for a device event (Section~\ref{sec:events}).

\subsection{Pending-access traps}
For debugging, setting bit 0 of CSR \csr{ccr} (PTRAP) turns the stall of a load, store, \insn{ltag} or
\insn{linst} on an in-flight range into a trap with cause \cause{pend} and \csr{tval} = the Local
Store address. The wait of a \insn{tget} or \insn{tput} whose D-LS range overlaps a pending device read
becomes a \cause{pend} trap in the same way, with \csr{tval} = the transfer's D-LS address \reg{rl}. This
shows where a program relies on the hardware stall. Whether a \cause{pend} trap occurs on a memory
transfer's range depends on timing; in the reference emulator, where memory transfers complete at
issue, it occurs only for the range of a pending device read. Instruction fetch stalls on an in-flight
\insn{tgetx}, and waits for a busy tag, are never converted into traps.

\section{Device transfers}\label{sec:devxfer}

Transfers to device space (space 4) differ from memory transfers.

\begin{rule*}{One byte per beat, to consecutive offsets}
A device transfer moves one byte per beat. Beat $i$ of a transfer of $n$ bytes at offset \reg{ro}
addresses device offset $\texttt{base} + \reg{ro} + i$. There is no alignment requirement on \reg{rl},
\reg{ro} or \reg{rn}.
\end{rule*}

\begin{itemize}
\item The offset names where the transfer \emph{starts}: the bounds check requires only
  $\reg{ro} < \textit{len}$, and the length is not bounded by the capability. A device capability
  therefore grants every transfer that starts inside it.
\item A \term{register} is 8 bytes wide, little-endian, at an 8-byte-aligned offset. It is read or
  written by an 8-byte transfer at its offset. Its value is \emph{latched when byte 0 is read}, so a read
  never sees a value that changed between beats, and a written value \emph{takes effect when byte 7 is
  written}. A longer transfer runs on into the following registers: 16 bytes read from the timer's FREQ
  register return FREQ and then DEADLINE.
\item A \term{window} (the platform table, the block device's identify data) is read-only and
  byte-addressed like memory: a read of 4 bytes at offset 4 of the platform table returns its bytes
  4--7.
\item A \term{stream port} (console output, console input, the block device's submission and completion
  ports) takes the whole transfer that starts at it: every beat of a \insn{tput} that starts at the
  console output port is the next byte of output, and every beat of a \insn{tget} that starts at a stream
  port is the next byte of the stream.
\item Reads of unassigned offsets, and of write-only ports, return zero; writes to unassigned offsets
  and to read-only registers and windows are ignored.
\item \insn{tgetx} from device space is not permitted (\cause{perm}).
\end{itemize}

\subsection{Pending device reads}\label{sec:pendread}

A read that starts at an \term{event port} completes only when the device has something to deliver.
The prototype has four: console input (the next byte), the console break port (Ctrl-C), the timer's
WAIT register (the deadline) and the framebuffer controller's VSYNC register (the next refresh). Such a
read is a \term{pending device read}. It is an ordinary transfer in every respect but its latency:

\begin{itemize}
\item its tag is busy until the device completes it, which \insn{tpoll} reports and \insn{tany} and
  \insn{twait} wait for;
\item nothing else is held up: the core runs on, other cores run, the timer counts and the display
  refreshes;
\item a load or store that overlaps its D-LS destination waits until it completes (\cause{pend} with
  CCR.PTRAP);
\item a \insn{tget} or \insn{tput} whose D-LS range overlaps its destination waits \emph{at issue} until
  it completes, so transfers stay ordered after the read. With CCR.PTRAP it traps with \cause{pend}
  instead, with \csr{tval} = its D-LS address. \insn{tgetx} writes the I-LS and is not affected;
\item a core has at most four pending device reads; issuing a fifth waits. Software must keep at most
  one read of the break port in flight per core;
\item \insn{halt} drains them like every other transfer of the core, in supervisor mode before the core
  stops and in user mode before the \cause{halt} trap. Software must not halt a core while it has a
  read in flight that may never complete: a break read is retired first by writing the console's
  release port (Section~\ref{sec:break}); a WAIT is retired by writing a DEADLINE of zero.
\end{itemize}
Pending device reads do not hold up hand-overs: \insn{cgive}, \insn{csend}, \insn{cshare},
\insn{catomic} and \insn{tamo} wait only for the core's \emph{memory} transfers
(Section~\ref{sec:publication}).

The ports, registers and windows of the prototype platform are defined in Chapter~\ref{ch:platform}.

\section{Far atomics}\label{sec:atomics}

\insn{tamo.add}, \insn{tamo.swap}, \insn{tamo.cas}, \insn{tamo.and} and \insn{tamo.or} perform an
indivisible read-modify-write of one aligned doubleword at the memory that holds it, and return the old
value. They require the \perm{A} permission, which only \insn{catomic} grants; an atomic capability has
no \perm{R} or \perm{W}, so an atomic segment is reachable \emph{only} through atomics. All atomics on
one doubleword, from all cores, are linearizable.

A far atomic first waits until none of the issuing core's memory transfers is in flight, so every
far access the core issued before it is complete and visible to every other core before the atomic is (the
publication rule, Section~\ref{sec:publication}). Atomics do not use a tag and are complete when the
instruction retires.

\section{Timing-visible state}\label{sec:timing-visible}

Apart from the following, nothing a single core can observe depends on transfer latency or on the
timing of other cores:
\begin{itemize}
\item the \csr{cycle} and \csr{stalls} CSRs;
\item the results of \insn{tpoll} and \insn{tany};
\item \cause{pend} traps when CCR.PTRAP is set.
\end{itemize}
Inter-core interactions, stop requests and device events add a few more sources of timing
dependence, listed in Section~\ref{sec:determinism}.
